\begin{table*}[t]
\centering
\scriptsize
\renewcommand{\arraystretch}{0.9}
\begin{tabularx}{\linewidth}{>{\raggedright\arraybackslash}p{0.20\linewidth}>{\raggedright\arraybackslash}p{0.24\linewidth}>{\raggedright\arraybackslash}p{0.20\linewidth}X}
\toprule
\textbf{Recall check} & \textbf{Observed result} & \textbf{What it controls} & \textbf{Residual risk} \\
\midrule
ACL Anthology metadata scan & 12,924 entries scanned for ACL/EMNLP/NAACL/COLING; 181 RAG candidates; 30 robust-RAG candidates & Reduces reliance on arXiv-only discovery for core NLP venues. & Title/abstract terms can miss papers using older QA/IR terminology. \\
Cross-venue metadata pass & DBLP/OpenAlex queries yielded 51 venue-linked robust-RAG candidates & Adds IR, web, security, and ML venue-linked records. & API ranking, metadata naming, and transient query errors limit recall. \\
Follow-up triage & 51 candidates triaged; 26 added to scoping set; 24 added to evidence map; 9 excluded; 2 duplicates & Closes the known unprocessed venue-candidate gap. & Does not prove complete coverage of every proceedings volume. \\
Evidence-strength sensitivity & Main conclusions separated into venue-linked backbone vs preprint trend evidence & Prevents preprint-heavy map statistics from becoming field-level claims. & Some emerging security/control subtopics remain too new for strong consensus. \\
\bottomrule
\end{tabularx}
\caption{Venue-recall risk analysis. The venue pass improves auditability and adds missed venue-linked papers, but it is still not a manual venue-complete proceedings review.}
\label{tab:venue-recall-risk}
\end{table*}

